\section*{Приложение A. Исходный код всех программ, использовавшихся в лабораторной работе (в комментариях приведены названия файлов)}
\addcontentsline{toc}{section}{Приложение A. Исходный код всех программ, использовавшихся в лабораторной работе (в комментариях приведены названия файлов)}

\renewcommand{\theFancyVerbLine}{%
    \scriptsize\ttfamily\textcolor{gray}{\arabic{FancyVerbLine}}%
}

\begin{minted}[linenos, numbersep=10pt, firstnumber=1, fontsize=\scriptsize]{matlab}
% ** task1a.m **

f = @(x) abs(x .* sin(x) - 4 .* cos(x));

[x_max, neg_max] = fminbnd(@(x) -f(x), 0, pi);
M4 = -neg_max;

fprintf('x_max = %.4f\n', x_max);
fprintf('M4 = %.4f\n', M4);

% ** task1b.m **
f = @(x) x .* sin(x);

a = 0;
b = pi;

h = (b - a) / 3;
x = [a, a + h, a + 2 * h, b];
y = f(x);

I_nc = 3 * h / 8 * (y(1) + 3 * y(2) + 3 * y(3) + y(4));
I_exact = pi;
error_nc = abs(I_exact - I_nc);

fprintf('I_nc = %.6f\n', I_nc);
fprintf('error_nc = %.6f\n', error_nc);

% ** task2.m **
f = @(x) x .* sin(x);

a = 0;
b = pi;

nodes = [-0.9061798459, -0.5384693101, 0, ...
          0.5384693101,  0.9061798459];
weights = [0.2369268851, 0.4786286705, 0.5688888889, ...
           0.4786286705, 0.2369268851];

x = (a + b) / 2 + (b - a) / 2 * nodes;
y = f(x);

I_gauss = (b - a) / 2 * sum(weights .* y);
I_exact = pi;
error_gauss = abs(I_exact - I_gauss);

fprintf('I_gauss = %.12f\n', I_gauss);
fprintf('error_gauss = %.12f\n', error_gauss);
\end{minted}
